% EDITION STATUS: partial; source units 21/116
\documentclass[11pt,a4paper]{article}
\usepackage[no-math]{fontspec}
\setmainfont{Latin Modern Roman}
\setsansfont{Latin Modern Sans}
\setmonofont{Latin Modern Mono}
\usepackage[english,french]{babel}
\usepackage[left=22mm,right=34mm,top=22mm,bottom=22mm,marginparwidth=23mm,marginparsep=5mm,includeheadfoot]{geometry}
\usepackage{amsmath,amssymb}
\usepackage[normalem]{ulem}
\usepackage{xcolor,tikz-cd,marginnote,etoolbox,microtype}
\usepackage[unicode,colorlinks=true,linkcolor=LinkInk,urlcolor=LinkInk,bookmarksopen=true]{hyperref}
\definecolor{LinkInk}{HTML}{254E70}
\definecolor{HandInk}{HTML}{254E70}
\definecolor{DeletionInk}{HTML}{8B3A3A}
\definecolor{QuietInk}{HTML}{555555}
\setlength{\parindent}{1em}
\setlength{\parskip}{3pt plus 1pt}
\setlength{\emergencystretch}{1.5em}
\raggedbottom
\normalmarginpar
\frenchspacing
\newif\ifdiplomatic
\newif\ifreading
\newif\ifenglish
\diplomatictrue
\readingfalse
\englishfalse
\newcommand{\EditionTitle}{Esquisse d'un programme}
\newcommand{\EditionForm}{Texte diplomatique — les deux copies}
\hypersetup{pdftitle={Esquisse d'un programme — Texte diplomatique — les deux copies},pdfauthor={Alexander Grothendieck},pdfsubject={A scan-checked edition of Esquisse d'un programme, January 1984, Montpellier cote 119}}
\newcommand{\DeclareScan}[4]{%
  \expandafter\def\csname edurl@#1\endcsname{#2}%
  \expandafter\def\csname edkey@#1\endcsname{#3}%
  \expandafter\def\csname edpage@#1\endcsname{#4}}
\def\CurrentScanURL{}
\def\CurrentPageKey{}
\def\CurrentCopy{copy1}
\def\CurrentApparatus{}
\newcounter{appnote}
\newcommand{\EditionPart}[2]{\clearpage\def\CurrentCopy{#1}\section*{#2}}
\newcommand{\scanpage}[1]{%
  \clearpage
  \expandafter\let\expandafter\CurrentScanURL\csname edurl@#1\endcsname
  \expandafter\let\expandafter\CurrentPageKey\csname edkey@#1\endcsname
  \def\CurrentApparatus{}%
  \hypertarget{scan-\CurrentPageKey}{}%
  \typeout{EDITION-SCAN: \CurrentPageKey; PDFPAGE: \the\numexpr\ReadonlyShipoutCounter+1\relax}%
  \ifreading
    \marginnote{\raggedright\scriptsize\sffamily\href{\CurrentScanURL}{119 / \csname edpage@#1\endcsname}}%
  \else
    \noindent{\small\sffamily\href{\CurrentScanURL}{[scan #1]}}\par\smallskip
  \fi}
\newcommand{\mspage}[1]{%
  \ifreading\marginnote{\raggedright\scriptsize\sffamily p.\,#1}[2\baselineskip]%
  \else\noindent{\small\sffamily[p.\,#1]}\par\smallskip\fi}
\newcommand{\hand}[1]{\ifdiplomatic\textcolor{HandInk}{#1}\else#1\fi}
\let\over\relax
\newcommand{\over}[2]{\textcolor{QuietInk}{#1}\textcolor{HandInk}{\ensuremath{^{\!\nearrow}}#2}}
\newcommand{\inserted}[2][]{\textcolor{HandInk}{\ensuremath{\langle}#2\ensuremath{\rangle}}}
\newcommand{\deleted}[1]{\ifmmode\textcolor{DeletionInk}{[#1]_{\scriptscriptstyle\mathrm{del}}}\else\textcolor{DeletionInk}{\sout{#1}}\fi}
\renewcommand{\xout}[1]{\ifmmode\textcolor{QuietInk}{[#1]_{\scriptscriptstyle\mathrm{x}}}\else\textcolor{QuietInk}{\sout{#1}}\fi}
\let\TextUline\uline
\let\TextUuline\uuline
\renewcommand{\uline}[1]{\ifmmode\underline{#1}\else\TextUline{#1}\fi}
\renewcommand{\uuline}[1]{\ifmmode\underline{\underline{#1}}\else\TextUuline{#1}\fi}
\newcommand{\unsure}[1]{\href{\CurrentScanURL}{[#1?]}}
\newcommand{\illegible}[1]{\href{\CurrentScanURL}{[\ifenglish illegible\else illisible\fi\ #1]}}
\newsavebox{\MarginalTestBox}
\newcommand{\margin}[1]{%
  \sbox{\MarginalTestBox}{\parbox{\marginparwidth}{\raggedright\footnotesize#1}}%
  \ifdim\dimexpr\ht\MarginalTestBox+\dp\MarginalTestBox\relax>3\baselineskip
    \par\smallskip\noindent\begin{minipage}{\linewidth}\hrule\smallskip
    {\footnotesize\textcolor{HandInk}{\ifenglish Marginal note:\else Note marginale :\fi\ #1}}\smallskip\hrule\end{minipage}\par\smallskip
  \else
    \marginnote{\raggedright\footnotesize\textcolor{HandInk}{#1}}%
  \fi}
\newcommand{\authornote}[1]{\margin{#1}}
\newcommand{\diagram}[1]{\par\smallskip{\centering#1\par}\smallskip}
\newenvironment{diagramblock}{\par\begin{center}}{\end{center}\par}
\newcommand{\DeclareVariant}[2]{\expandafter\def\csname edvariant@#1\endcsname{#2}}
\newcommand{\AppItem}[2]{\par\noindent\textsuperscript{#1}\csname edvariant@#2\endcsname}
\newcommand{\EditionApp}[1]{%
  \stepcounter{appnote}%
  \ifmmode{}^{\scriptscriptstyle\theappnote}\else\textsuperscript{\theappnote}\fi
  \begingroup\edef\x{\endgroup\noexpand\gappto\noexpand\CurrentApparatus{\noexpand\AppItem{\theappnote}{#1}}}\x}
\newcommand{\var}[4]{\ifdiplomatic#1\else#2\ifreading\else\EditionApp{#4}\fi\fi}
\newcommand{\copytwo}[2]{\ifdiplomatic\else\ifreading\else\DeclareVariant{#1}{#2}\EditionApp{#1}\fi\fi}
\newcommand{\tnote}[2]{\ifreading\else\DeclareVariant{#1}{\textit{Translator's note.} #2}\EditionApp{#1}\fi}
\newcommand{\finishscan}{\ifx\CurrentApparatus\empty\else\par\bigskip\noindent\rule{.35\linewidth}{.35pt}\par\smallskip{\footnotesize\setlength{\parskip}{3pt}\CurrentApparatus}\fi}
\newcommand{\DeclareNoteReturns}[2]{\expandafter\def\csname edreturns@#1\endcsname{#2}}
\newcommand{\notecall}[3]{\hypertarget{#2}{\hyperlink{note-#1}{#3}}}
\newcommand{\notedef}[2]{\hypertarget{note-#1}{#2}\,{\scriptsize[\csname edreturns@#1\endcsname]}}
\newcommand{\AppReading}[1]{{\def\margin##1{##1}\def\authornote##1{##1}#1}}
\newcommand{\EmptyReading}{\ifenglish\textit{omitted}\else\textit{supprimé}\fi}
\DeclareScan{119/2}{https://grothendieck.umontpellier.fr/web/viewer.html?file=../119.pdf\#page=3}{esq_119_002}{2}
\DeclareScan{119/3}{https://grothendieck.umontpellier.fr/web/viewer.html?file=../119.pdf\#page=4}{esq_119_003}{3}
\DeclareScan{119/4}{https://grothendieck.umontpellier.fr/web/viewer.html?file=../119.pdf\#page=5}{esq_119_004}{4}
\DeclareScan{119/5}{https://grothendieck.umontpellier.fr/web/viewer.html?file=../119.pdf\#page=6}{esq_119_005}{5}
\DeclareScan{119/6}{https://grothendieck.umontpellier.fr/web/viewer.html?file=../119.pdf\#page=7}{esq_119_006}{6}
\DeclareScan{119/7}{https://grothendieck.umontpellier.fr/web/viewer.html?file=../119.pdf\#page=8}{esq_119_007}{7}
\DeclareScan{119/8}{https://grothendieck.umontpellier.fr/web/viewer.html?file=../119.pdf\#page=9}{esq_119_008}{8}
\DeclareScan{119/9}{https://grothendieck.umontpellier.fr/web/viewer.html?file=../119.pdf\#page=10}{esq_119_009}{9}
\DeclareScan{119/10}{https://grothendieck.umontpellier.fr/web/viewer.html?file=../119.pdf\#page=11}{esq_119_010}{10}
\DeclareScan{119/11}{https://grothendieck.umontpellier.fr/web/viewer.html?file=../119.pdf\#page=12}{esq_119_011}{11}
\DeclareScan{119/12}{https://grothendieck.umontpellier.fr/web/viewer.html?file=../119.pdf\#page=13}{esq_119_012}{12}
\DeclareScan{119/13}{https://grothendieck.umontpellier.fr/web/viewer.html?file=../119.pdf\#page=14}{esq_119_013}{13}
\DeclareScan{119/14}{https://grothendieck.umontpellier.fr/web/viewer.html?file=../119.pdf\#page=15}{esq_119_014}{14}
\DeclareScan{119/15}{https://grothendieck.umontpellier.fr/web/viewer.html?file=../119.pdf\#page=16}{esq_119_015}{15}
\DeclareScan{119/16}{https://grothendieck.umontpellier.fr/web/viewer.html?file=../119.pdf\#page=17}{esq_119_016}{16}
\DeclareScan{119/17}{https://grothendieck.umontpellier.fr/web/viewer.html?file=../119.pdf\#page=18}{esq_119_017}{17}
\DeclareScan{119/18}{https://grothendieck.umontpellier.fr/web/viewer.html?file=../119.pdf\#page=19}{esq_119_018}{18}
\DeclareScan{119/19}{https://grothendieck.umontpellier.fr/web/viewer.html?file=../119.pdf\#page=20}{esq_119_019}{19}
\DeclareScan{119/20}{https://grothendieck.umontpellier.fr/web/viewer.html?file=../119.pdf\#page=21}{esq_119_020}{20}
\DeclareScan{119/21}{https://grothendieck.umontpellier.fr/web/viewer.html?file=../119.pdf\#page=22}{esq_119_021}{21}
\DeclareScan{119/22}{https://grothendieck.umontpellier.fr/web/viewer.html?file=../119.pdf\#page=23}{esq_119_022}{22}
\DeclareVariant{V-FR-006-001}{\textit{coquille corrigée}. \AppReading{des surface.} \(\rightarrow\) \AppReading{des surfaces.}. Le pluriel est imposé localement par « des »; le même substantif est écrit « surfaces » à la page 3. La correction ajoute le s final manquant.}
\DeclareVariant{V-FR-007-001}{\textit{coquille corrigée}. \AppReading{voir indispensable} \(\rightarrow\) \AppReading{voire indispensable}. La progression « bien plus élégant, voire indispensable » appelle l'adverbe d'enchérissement; la forme « voire » figure à la page 2 (« voire conformes ») et à la page 3. Le e final a été omis ici.}
\DeclareVariant{V-FR-009-001}{\textit{coquille corrigée}. \AppReading{convenanble} \(\rightarrow\) \AppReading{convenable}. Un n supplémentaire est intercalé dans le mot; « convenable » figure dans le même développement à la page 5.}
\DeclareVariant{V-FR-010-001}{\textit{coquille corrigée}. \AppReading{Ladegraillerie} \(\rightarrow\) \AppReading{Ladegaillerie}. Le nom « Yves Ladegaillerie » figure à la ligne précédente pour le même participant; le second emploi intercale un r.}

\begin{document}
\begin{titlepage}
\centering
\vspace*{18mm}
{\Large Alexander Grothendieck\par}
\vspace{12mm}
{\Huge\bfseries\EditionTitle\par}
\vspace{8mm}
{\Large\EditionForm\par}
\vspace{14mm}
{\large January 1984\par}
\vfill
{\large Université de Montpellier\par fonds Grothendieck, cote 119\par}
\vspace{8mm}
{\normalsize Two typescript copies: sites 2--59 and 78--135\par}
\vspace{8mm}
{\normalsize Edition dated 14 November 2026\par}
\end{titlepage}
\thispagestyle{plain}
\section*{Préface}
\noindent\textbf{Édition de travail partielle.} Ce fichier contient 21 pages sources sur les 116 prévues pour cette forme. La collation et les contrôles finaux se poursuivent.\par\smallskip
Ce volume donne intégralement la copie 1 puis la copie 2. Chaque page porte son lien vers le fac-similé de la cote 119. Le bleu marque l’écriture manuscrite ; le texte biffé, les insertions, les surcharges et les notes marginales conservent leur distinction. Les mots soulignés restent soulignés. Les lectures incertaines restent signalées par des crochets et un point d’interrogation, avec un lien vers le scan. Les numéros de page du document et les marques d’archive sont conservés. Les notes (1) à (7) restent à la fin, avec des liens de retour.

\smallskip
AI-transcribed by Claude Opus 5.5 (Anthropic, effort high, Message Batches API); checked against the scans, and this edition built, by GPT-6 Astra (OpenAI, ChatGPT Pro mode); prompts and records by Claude (Anthropic).


The transcription of copy 1 was made with the text layer of the Grothendieck Circle's typeset \textit{Esquisse} (\texttt{EsquisseFr.pdf}) given to the transcribing model as a draft, to compare with after it had transcribed each page; every reading of this edition was decided on the scans.

The published French text and English translation in \textit{Geometric Galois Actions 1}, edited by L. Schneps and P. Lochak (Cambridge University Press, 1997), the Grothendieck Circle's typeset \texttt{EsquisseFr.pdf}, and the Centre for Grothendieckian Studies' draft transcription G94d are named as published texts. They are not reproduced here.

\section*{Droits et crédits}
The rights in the original work belong to its author or the author's representatives or heirs. Where we hold rights, in the AI-produced transcription, translation, corrections and typesetting, they are dedicated to the public domain under CC0 1.0. This edition is published for educational, non-profit purposes, and we claim no rights we do not hold. If you hold rights in this work and object to its publication here, please open a GitHub issue and we will take it down.

Les droits sur l'œuvre originale appartiennent à son auteur, à ses représentants ou à ses ayants droit. Dans la mesure où nous détenons des droits sur la transcription, la traduction, les corrections et la composition produites par l'IA, nous les consacrons au domaine public sous CC0 1.0. Cette édition est publiée à des fins éducatives et non lucratives, et nous ne revendiquons aucun droit que nous ne détenons pas. Si vous détenez des droits sur cette œuvre et vous opposez à sa publication ici, veuillez ouvrir une issue GitHub et nous la retirerons.
\clearpage

\EditionPart{copy1}{Copie 1}
% BEGIN PAGE esq_119_002
\scanpage{119/2}
\uline{ESQUISSE D'UN PROGRAMME}

par Alexandre Grothendieck

\uline{Sommaire} :

1. Envoi.

2. Un "jeu de Légo-Teichmüller", et le groupe de Galois de $\overline{\mathbb{Q}}$ sur $\mathbb{Q}$.

3. Corps de nombres associés à un dessin d'enfant.

4. Polyèdres réguliers et corps finis.

5. Haro sur la topologie dite ``générale'', et réflexions heuristiques
vers une topologie dite ``modérée''.

6. ``Théories différentiables'' (à la Nash) et ``théories modérées''.

7. A la Poursuite des Champs.

8. Digressions de géométrie bidimensionelle.

9. Bilan d'une activité enseignante.

10. Epilogue.

Notes.

N.B. Les astéristiques (*) renvoient aux notes figurant au bas de
la même page, les renvois numérotés de (\textsuperscript{\hand{1}}) à (\textsuperscript{\hand{7}}\unsure{.}) aux notes (rajou-%
tées ultérieurement) réunies à la fin du rapport.

\mspage{\hand{2}}

\finishscan
% END PAGE esq_119_002
% BEGIN PAGE esq_119_003
\scanpage{119/3}
\mspage{1}

ESQUISSE D'UN PROGRAMME

par A. Grothendieck

1. Comme la conjoncture actuelle rend de plus en plus illu-%
soires pour moi les perspectives d'un enseignement de recherche à
l'Université, je me suis résolu à demander mon admission au CNRS, pour
pouvoir consacrer mon énergie à développer des travaux et perspectives
dont il devient clair qu'il ne se trouvera aucun élève (ni même,
semble-t-il, aucun congénère mathématicien) pour les développer à ma
place.

En guise de document ``Titres et Travaux'', on trouvera à la
suite de ce texte la reproduction intégrale d'une esquisse, par thèmes,
de ce que je considérais comme mes principales contributions mathéma-%
tiques au moment d'écrire ce rapport, en 1972. Il contient également
une liste d'articles publiés à cette date. J'ai cessé toute publica-%
tion d'articles scientifiques depuis 1970. Dans les lignes qui suivent,
je me propose de donner un aperçu au moins sur quelques thèmes princi-%
paux de mes réflexions mathématiques depuis lors. Ces réflexions se
sont matérialisées au cours des années en deux volumineux cartons de
notes manuscrites, difficilement déchiffrables sans doute à tout autre
qu'à moi-même, et qui, après des stades de décantations successives,
attendent leur heure peut-être pour une rédaction d'ensemble tout au
moins provisoire, à l'intention de la communauté mathématique. Le
terme ``rédaction'' ici est quelque peu impropre, alors qu'il s'agit
bien plus de développer des idées et visions multiples amorcées au
cours de ces douze dernières années, en les précisant et les appro-%
fondissant, avec tous les rebondissements imprévus qui constamment
accompagnent ce genre de travail - un travail de découverte donc, et
non de compilation de notes pieusement accumulées. Et je compte bien,
dans l'écriture des ``Réflexions Mathématiques'' commencée depuis fé-%
vrier 1983, laisser apparaître clairement au fil des pages la démarche
de la pensée qui sonde et qui découvre, en tâtonnant dans la pénombre
bien souvent, avec des trouées de lumière subites quand quelque tenace

\margin{3}

\finishscan
% END PAGE esq_119_003
% BEGIN PAGE esq_119_004
\scanpage{119/4}
\mspage{2}

image fausse, ou simplement inadéquante, se trouve enfin débusquée
et mise à jour, et que les choses qui semblaient de guingois se
mettent en place, dans l'harmonie mutuelle qui leur est propre.
\qquad Quoi qu'il en soit, l'esquisse qui suit de quelques
thèmes de réflexion des dernières dix ou douze années, tiendra
lieu en même temps d'esquisse de programme de travail pour les
années qui viennent, que je compte consacrer au développement de
ces thèmes, ou au moins de certains d'entre eux. Elle est desti-%
née, d'une part aux collègues du Comité National appelés à statuer
sur ma demande, d'autre part à quelques autres collègues, anciens
élèves, amis, dans l'éventualité où certaines des idées esquissées
ici pourraient intéresser l'un d'entre eux.

2. \qquad Les exigences d'un enseignement universitaire, s'adressant
donc à des étudiants (y compris les étudiants dits ``avancés'') au ba-%
gage mathématique modeste (et souvent moins que modeste), m'ont amené
à renouveler de façon draconienne les thèmes de réflexion à proposer
à mes élèves, et de fil en aiguille et de plus en plus, à moi-même
également. Il m'avait semblé important de partir d'un bagage intuitif
commun, indépendant de tout langage technique censé l'exprimer, bien
antérieur à tout tel langage - il s'est avéré que l'intuition géomé-%
trique et topologique des formes, et plus particulièrement des formes
bidimensionelles, était un tel terrain commun. Il s'agit donc de thèmes
qu'on peut grouper sous l'appellation de ``topologie des surfaces'' ou
``géométrie des surfaces'', étant entendu dans cette dernière appella-%
tion que l'accent principal se trouve sur les propriétés topologiques
des surfaces, ou sur les aspects combinatoires qui en constituent
l'expression technique la plus terre-à-terre, et non sur les aspects
différentiels, voire conformes, riemaniens, holomorphes et (de là)
l'aspect ``courbes algébriques complexes''. Une fois ce dernier pas
franchi cependant, voici soudain la géométrie algébrique (mes anciennes
amours !) qui fait irruption à nouveau, et ce par les objets qu'on
peut considérer comme les pierres de construction ultimes de toutes
les autres variétés algébriques. Alors que dans mes recherches d'avant

\margin{4}

\finishscan
% END PAGE esq_119_004
% BEGIN PAGE esq_119_005
\scanpage{119/5}
\mspage{3}

1970, mon attention systématiquement était dirigée vers les objets
de généralité maximale, afin de dégager un langage d'ensemble adéquat
pour le monde de la géométrie algébrique, et que je ne m'attardais
sur les courbes algébriques que dans la stricte mesure où cela s'avé-%
rait indispensable (notamment en cohomologie étale) pour développer
des techniques et énoncés ``passe-partout'' valables en toutes dimen-%
sions et en tous lieux (j'entends, sur tous schémas de base, voire
tous topos annelés de base ...), me voici donc ramené, par le tru-%
chement d'objets si simples qu'un enfant peut les connaître en jouant,
aux débuts et origines de la géométrie algébrique, familiers à Riemann
et à ses émules !

Depuis environ 1975, c'est donc la géométrie des surfaces
(réelles), et à partir de 1977 les liens entre les questions de géo-%
métrie des surfaces et la géométrie algébrique des courbes algébri-%
ques définies sur des corps tels que $\mathbb{C}$, $\mathbb{R}$ ou des extensions de type
fini de $\mathbb{Q}$, qui ont été ma principale source d'inspiration, ainsi
que mon fil conducteur constant. C'est avec surprise et avec émer-%
veillement qu'au fil des ans je découvrais (ou plutôt, sans doute,
redécouvrais) la richesse prodigieuse, réellement inépuisable, la
profondeur insoupçonnée de ce thème, d'apparence si anodine. Je crois
y sentir un point névralgique entre tous, un point de convergence
privilégié des principaux courants d'idées mathématiques, comme aussi
des principales structures et des visions des choses qu'elles expri-%
ment, depuis les plus spécifiques (tels les anneaux $\mathbb{Z}$, $\mathbb{Q}$, $\bar{\mathbb{Q}}$, $\mathbb{R}$, $\mathbb{C}$,
ou le groupe $S\ell(2)$ sur l'un de ces anneaux, ou les groupes algébri-%
ques réductifs généraux) aux plus ``abstraits'', telles les ``multipli-%
cités'' algébriques, analytiques complexes ou analytiques réelles.
(Celles-ci s'introduisent naturellement quand il s'agit d'étudier
systématiquement des ``variétés de modules'' pour les objets géométri-%
ques envisagés, si on veut dépasser le point de vue notoirement in-%
suffisant des ``modules grossiers'', qui revient à ``tuer'' bien malen-%
contreusement les groupes d'automorphismes de ces objets .) Parmi ces
multiplicités modulaires, ce sont celles de Mumford-Deligne pour les
courbes algébriques ``stables'' de genre $g$, à $\nu$ points marqués, que

\margin{5}

\finishscan
% END PAGE esq_119_005
% BEGIN PAGE esq_119_006
\scanpage{119/6}
\mspage{4}

je note $M^{\wedge}_{g,\nu}$ (compactification de la multiplicité ``ouverte''
$M_{g,\nu}$ correspondant aux courbes lisses), qui depuis quelques deux
ou trois années ont exercé sur moi une fascination particulière,
plus forte peut-être qu'aucun autre objet mathématique à ce jour.
A vrai dire, il s'agit plutôt du système de \uline{toutes} les multiplicités
$M_{g,\nu}$ pour $g,\nu$ variables, liées entre elles par un certain nombre
d'opérations fondamentales (telles les opérations de ``bouchage de
trous'' i.e. de ``gommage'' de points marqués, celle de ``recollement'',
et les opérations inverses), qui sont le reflet en géométrie algé-%
brique absolue de caractéristique zéro (pour le moment) d'opérations
géométriques familières du point de vue de la ``chirurgie'' topologique
ou conforme \var{des surface.}{des surfaces.}{slip corrected}{V-FR-006-001} La principale raison sans doute de cette
fascination, c'est que cette structure géométrique très riche sur le
système des multiplicités modulaires ``ouvertes'' $M_{g,\nu}$ se réflète
par une structure analogue sur les groupoïdes fondamentaux corres-%
pondants, les ``groupoïdes de Teichmüller'' $\hat{T}_{g,\nu}$ ,et que ces opérations
au niveau des $\hat{T}_{g,\nu}$ ont un caractère suffisamment intrinsèque pour
que le groupe de Galois $\mathrm{I}\!\Gamma$ de $\overline{\mathbb{Q}}/\mathbb{Q}$ opère sur toute cette ``tour'' de
groupoïdes de Teichmüller, en respectant toutes ces structures. Chose
plus extraordinaire encore, cette opération est \uline{fidèle} - à vrai dire,
elle est fidèle déjà sur le premier ``étage'' non trivial de cette tour,
savoir $\hat{T}_{0,4}$ - ce qui signifie aussi, essentiellement, que l'action
extérieure de $\mathrm{I}\!\Gamma$ sur le groupe fondamental $\hat{\pi}_{0,3}$ de la droite pro-%
jective standard $\mathbb{P}^1$ sur $\mathbb{Q}$, privée des trois points $0,\ 1,\ \infty$, est
déjà fidèle. Ainsi \uline{le groupe de Galois} $\mathrm{I}\!\Gamma$ \uline{se réalise comme un groupe}
\uline{d'automorphisme d'un groupe profini des plus concrets}, respectant
d'ailleurs certaines structures essentielles de ce groupe. Il s'ensuit
qu'un élément de $\mathrm{I}\!\Gamma$ peut être ``paramétré'' (de diverses façons équi-%
valentes d'ailleurs) par un élément convenable de ce groupe profini
$\hat{\pi}_{0,3}$ (un groupe profini libre à deux générateurs), ou par un système
de tels éléments, ce ou ces éléments étant d'ailleurs soumis à cer-%
taines conditions simples, nécessaires (et sans doute non suffisantes)
pour que ce ou ces éléments corresponde(nt) bien à un élément de $\mathrm{I}\!\Gamma$.
Une des tâches les plus fascinantes ici, est justement d'appréhender
des conditions nécessaires \uline{et} suffisantes sur un automorphisme extérieur

\margin{6}

\finishscan
% END PAGE esq_119_006
% BEGIN PAGE esq_119_007
\scanpage{119/7}
\mspage{5}

de $\hat{\pi}_{0,3}$ i.e. sur le ou les paramètres correspondants, pour qu'il
provienne d'un élément de $\mathrm{I}\!\Gamma$ - ce qui fournirait une description
``purement algébrique'', en termes de groupes profinis et sans ré-%
férence à la théorie de Galois des corps de nombres, du groupe de
Galois $\mathrm{I}\!\Gamma = \mathrm{Gal}(\bar{\mathbb{Q}}/\mathbb{Q})$ !

Peut-être une caractérisation même conjecturale de $\mathrm{I}\!\Gamma$
comme sous-groupe de Autext($\hat{\pi}_{0,3}$) est-elle pour le moment hors de
portée\notecall{copy1-1}{call-esq_119_007-1-1}{($^1$)} ; je n'ai pas de conjecture à proposer encore. Une autre
tâche par contre est abordable immédiatement, c'est celle de décrire
l'action de $\mathrm{I}\!\Gamma$ sur toute la tour de Teichmüller, en termes de son
action sur le ``premier étage'' $\hat{\pi}_{0,3}$ , i.e. exprimer un automorphisme
de cette tour, en termes du ``paramètre'' dans $\hat{\pi}_{0,3}$ qui repère l'élé-%
ment courant $\gamma$ de $\mathrm{I}\!\Gamma$. Ceci est lié à une représentation de la tour
de Teichmüller (en tant que groupoïde muni d'une opération de ``recol-%
lement'') par générateurs et relations, qui donnera en particulier
une présentation par générateurs et relations, au sens ordinaire, de
chacun des $\hat{T}_{g,\nu}$ (en tant que groupoïde profini). Ici, même pour
$g = 0$ (donc quand les groupes de Teichmüller correspondants sont
des groupes de tresses ``bien connus''), les générateurs et relations
connus à ce jour dont j'ai eu connaissance, me semblent inutilisables
tels quels, car ils ne présentent pas les caractères d'invariance et
de symétrie indispensables pour que l'action de $\mathrm{I}\!\Gamma$ soit directement
lisible sur cette présentation. Ceci est lié notamment au fait que
les gens s'obstinent encore, en calculant avec des groupes fondamen-%
taux, à fixer un seul point base, plutôt que d'en choisir astucieu-%
sement tout un paquet qui soit invariant par les symétries de la si-%
tuation, lesquelles sont donc perdues en route. Dans certaines situa-%
tions (comme des théorèmes de descente à la Van Kampen pour groupes
fondamentaux) il est bien plus élégant, \var{voir indispensable}{voire indispensable}{slip corrected}{V-FR-007-001} pour y com-%
prendre quelque chose, de travailler avec des groupoïdes fondamentaux
par rapport à un paquet de points bases convenable , et il en est cer-%
tainement ainsi pour la tour de Teichmüller. Il semblerait (incroyable,
mais vrai !) que la géométrie même du premier étage de la tour de
Teichmüller (correspondant donc aux ``modules'' soit pour des droites

\margin{7}

\finishscan
% END PAGE esq_119_007
% BEGIN PAGE esq_119_008
\scanpage{119/8}
\mspage{6}

projectives avec quatre points marqués, soit pour des courbes ellip-%
tiques (!)) n'ait jamais été bien explicité, par exemple la relation
entre le cas de genre 0 avec la géométrie de l'octaèdre, et celle du
tétraèdre. A fortiori les multiplicités modulaires $M_{0,5}$ (pour les
droites projectives avec cinq points marqués) et $M_{1,2}$ (pour les
courbes de genre 1 avec deux points marqués), d'ailleurs quasiment
isomorphes entre elles, semblent-elles terre vierge - les groupes
de tresses ne vont pas nous éclairer à leur sujet ! J'ai commencé à
regarder $M_{0,5}$ à des moments perdus, c'est un véritable joyau, d'une
géométrie très riche étroitement liée à celle de l'icosaèdre.

L'intérêt a priori d'une connaissance complète des deux
premiers étages de la tour (savoir, les cas où la dimension modulaire
$N = 3g-3+\nu$ est $\leq 2$) réside dans ce principe, que \uline{la tour entière se
reconstitue à partir des deux premiers étages}, en ce sens que via
l'opération fondamentale de ``recollement'', l'étage 1 fournit un sys-%
tème complet de générateurs, et l'étage 2 un système complet de re-%
lations. Il y a une analogie frappante, et j'en suis persuadé, pas
seulement formelle, entre ce principe, et le principe analogue de
Demazure pour la structure des groupes algébriques réductifs, si on
remplace le terme ``étage'' ou ``dimension modulaire'' par ``rang semi-%
simple du groupe réductif''. Le lien devient plus frappant encore,
si on se rappelle que le groupe de Teichmüller $T_{1,1}$ (dans le contexte
discret transcendant maintenant, et non dans le contexte algébrique
profini, où on trouve les complétions profinies des premiers) n'est
autre que $S\ell(2,\mathbb{Z})$, i.e. le groupe des points entiers du schéma en
groupes simple de rang 1 ``absolu'' $S\ell(2)_{\mathbb{Z}}$. Ainsi, \uline{la pierre de cons-%
truction fondamentale pour la tour de Teichmüller}, \uline{est essentiellement
la même que celle pour ``la tour'' des groupes réductifs de tous rangs} -
un groupe d'ailleurs dont on peut dire sans doute qu'il est présent
dans toutes les disciplines essentielles des mathématiques.

Ce principe de construction de la tour de Teichmüller n'est
pas démontré à l'heure actuelle - mais je n'ai aucun doute qu'il ne
soit valable. Il résulterait (via une théorie de dévissage des struc-%
tures stratifiées - en l'occurence les $M^{\wedge}_{g,\nu}$ - qui resterait à écrire,
cf. par. 5) d'une propriété extrêmement plausible des multiplicités

\margin{8}

\finishscan
% END PAGE esq_119_008
% BEGIN PAGE esq_119_009
\scanpage{119/9}
\mspage{7}

\unsure{mcdulaires|modulaires} ouvertes $M_{g,\nu}$ dans le contexte analytique complexe,
à savoir que pour une dimension modulaire $N \geq 3$, le groupe fonda-%
mental de $M_{g,\nu}$ (i.e. le groupe de Teichmüller habituel $T_{g,\nu}$)
est isomorphe au ``groupe fondamental à l'infini'' i.e. celui d'un
``voisinage tubulaire de l'infini''. C'est là une chose bien familière
(due à Lefschetz essentiellement) pour une variété lisse \uline{affine} de
dimension $N \geq 3$. Il est vrai que les multiplicités modulaires ne
sont pas affines (sauf pour des petites valeurs de $g$), mais il suf-%
firait qu'une telle $M_{g,\nu}$ de dimension $N$ (ou plutôt, un revête-%
ment fini \var{convenanble}{convenable}{slip corrected}{V-FR-009-001}) soit réunion de $N$-2 ouverts affines, donc
que $M_{g,\nu}$ ne soit pas ``trop proche d'une variété compacte''.

N'ayant aucun doute sur ce principe de construction de la
tour de Teichmüller, je préfère laisser aux experts de la théorie
transcendante, mieux outillés que moi, le soin de prouver le néces-%
saire (s'il s'en trouve qui soit intéressé), pour expliciter plutôt,
avec tout le soin qu'elle mérite, la structure qui en découle pour
la tour de Teichmüller par générateurs et relations, dans le cadre
discret cette fois et non profini - ce qui revient, essentiellement,
à une compréhension complète des quatre multiplicités modulaires
$M_{0,4}$, $M_{1,1}$, $M_{0,5}$, $M_{1,2}$, et de leurs groupoïdes fondamentaux par rap-%
port à des ``points base'' convenablement choisis. Ceux-ci s'offrent
tout naturellement, comme les courbes algébriques complexes du type
$(g,\nu)$ envisagé, qui ont un groupe d'automorphismes (nécessairement
fini) plus grand que dans le cas générique (*). En y incluant la
sphère holomorphe à trois points marqués (provenant de $M_{0,3}$ i.e.

\rule{4cm}{0.4pt}

(*) Il faut y ajouter de plus les ``points-base'' provenant par opéra-%
tions de recollement de ``pièces'' du même type en dimension modulaire
inférieure. D'autre part, en dimension modulaire 2 (cas de $M_{0,5}$ et $M_{1,2}$),
il convient d'exclure les points de certaines familles à un paramètre
de courbes admettant un automorphisme exceptionnel d'ordre 2. Ces fa-%
milles constituent d'ailleurs sur les multiplicités envisagées des cour-%
bes rationnelles remarquables, qui me paraissent un ingrédient important
de la structure de ces multiplicités.

\margin{9}

\finishscan
% END PAGE esq_119_009
% BEGIN PAGE esq_119_010
\scanpage{119/10}
\mspage{8}

de l'étage O), on trouve \uline{douze} ``\uline{pièces de construction}'' \uline{fondamentales}
(6 de genre O, 6 de genre 1) dans un ``jeu de Légo-Teichmüller (grande
boite), où les points marqués sur les surfaces envisagées sont rem-%
placés par des ``trous'' à bord, de façon à avoir des surfaces à bord,
donc des pièces de construction qui peuvent s'assembler par frotte-%
ment doux comme dans le jeu de Légo ordinaire cher à nos enfants
(ou petits-enfants...). Par assemblage on trouve un moyen tout ce
qu'il y a de visuel pour construire tout type de surface (ce sont
ces assemblages essentiellement qui seront les ``points bases'' pour
notre fameuse tour), et aussi de visualiser les ``\uline{chemins}'' \uline{élémentaires}
par des opérations tout aussi concrètes tels des ``twists'', ou des
automorphismes des pièces du jeu, et d'écrire les \uline{relations fondamentales}
entre chemins composés. Suivant la taille (et le prix !) de la boite
de construction utilisée, on \unsure{trcuve|trouve} d'ailleurs de nombreuses descrip-%
tions différentes de la tour de Teichmüller par générateurs et rela-%
tions. La boite la plus petite est réduite à des pièces toutes iden-%
tiques, de type (O,3) - ce sont les ``pantalons'' de Thurston, et le
jeu de Légo-Teichmüller que j'essaie de décrire, issu de motivations
et de réflexions de géométrie algébrique absolue sur le corps $\mathbb{Q}$,
est très proche du jeu de ``chirurgie géodésique hyperbolique'' de
Thurston, dont j'ai appris l'existence l'an dernier par Yves Ladegaillerie.
Dans un microséminaire avec Carlos Contou-Carrère et Yves \var{Ladegraillerie}{Ladegaillerie}{slip corrected}{V-FR-010-001},
nous avons amorcé une réflexion dont un des objets est de confronter
les deux points de vue, qui se complètent mutuellement.

J'ajoute que chacune des douze pièces de construction de
la ``grande boite'' se trouve muni d'une décomposition cellulaire cano-%
nique, stable par toutes les symétries, ayant comme seuls sommets les
``points marqués'' (ou centres des trous), et comme arètes certains che-%
mins géodésiques (pour la structure riemanienne canonique sur la sphère
ou le tore envisagé) entre certaines paires de sommets (savoir ceux
qui se trouvent sur un même ``lieu réel'', pour une structure réelle
convenable de la courbe algébrique complexe envisagée). Par suite,
toutes les surfaces obtenues dans ce jeu par assemblage sont munies
de structures cellulaires canoniques, qui à leur tour (cf.§3 plus bas)

\margin{10}

\finishscan
% END PAGE esq_119_010
% BEGIN PAGE esq_119_011
\scanpage{119/11}
\mspage{9}

permettent de considérer ces surfaces comme associées à des courbes
algébriques complexes (et même sur $\overline{\mathbb{Q}}$) canoniquement déterminées.
Il y a là un jeu de chassé-croisé typique entre le combinatoire,
et l'algébrique complexe (ou mieux, l'algébrique sur $\overline{\mathbb{Q}}$).

La ``petite boite'' aux pièces toutes identiques, qui a le
charme de l'économie, donnera sans doute une description relative-%
ment compliquée pour les relations (compliquée, mais nullement inex-%
tricable !). La grande boite donnera lieu à des relations plus nom-%
breuses (du fait qu'il y a beaucoup plus de points-bases et de che-%
mins remarquables entre eux), mais à structure plus transparente. Je
prévois qu'en dimension modulaire 2, tout comme dans le cas plus ou
moins familier de la dimension modulaire 1 (avec notamment la des-%
cription de $S\ell(2,\mathbb{Z})$ par $(\rho,\sigma \mid \rho^3 = \sigma^2,\ \rho^4 = \sigma^6 = 1)$), on trouvera
un engendrement par les groupes d'automorphismes des trois types de
pièces pertinentes, avec des relations simples que je n'ai pas déga-%
gé à l'heure d'écrire ces lignes. Peut-être même trouvera-t-on un
principe de ce genre pour tous les $T_{g,\nu}$, ainsi qu'une décomposition
cellulaire de $M^{\wedge}_{g,\nu}$ généralisant celles qui se présentent sponta-%
nément pour $M^{\wedge}_{0,4}$ et $M^{\wedge}_{1,1}$, et que j'entrevois dès à présent pour
la dimension modulaire 2, en utilisant les hypersurfaces correspon-%
dant aux diverses \uline{structures réelles} sur les structures complexes
envisagées, pour effectuer le découpage cellulaire voulu.


3. Plutôt que de suivre (comme prévu) un ordre thématique ri-%
goureux, je me suis laissé emporter par ma prédilection pour un thème
particulièrement riche et brûlant, auquel je compte me consacrer
d'ailleurs prioritairement pendant quelque temps, à partir de la ren-%
trée 84/85. Je reprends donc l'exposé thématique là où je l'ai laissé,
tout au début du paragraphe précédent.

Mon intérêt pour les surfaces topologiques commence à poindre
en 1974, où je propose à Yves Ladegaillerie le thème de l'étude iso-%
topique des plongements d'un 1-complexe topologique dans une surface
compacte. Dans les deux années qui suivent, cette étude le conduit à
un remarquable théorème d'isotopie, donnant une description algébrique

\margin{11}

\finishscan
% END PAGE esq_119_011
% BEGIN PAGE esq_119_012
\scanpage{119/12}
\mspage{10}

complète des classes d'isotopie de plongements de tels 1-complexes,
ou de surfaces compactes à bord, dans une surface compacte orientée,
en termes de certains invariants combinatoires très simples, et des
groupes fondamentaux des protagonistes. Ce théorème, qui doit pou-%
voir s'étendre sans mal aux plongements d'un espace compact quelcon-%
que (triangulable pour simplifier) dans une surface compacte orientée,
redonne comme corollaires faciles plusieurs résultats classiques
profonds de théorie des surfaces, et notamment le théorème d'isoto-%
pie de Baer. Il va finalement être publié, séparément du reste (et
dix ans après, vu la dureté des temps...), dans Topology. Dans le
travail de Ladegaillerie figure également une description purement
algébrique, en termes de groupoïdes fondamentaux, de la catégorie
``isotopique'' des surfaces compactes $X$, munies d'un 1-complexe to-%
pologique $K$ plongé dans $X$. Cette description, qui a eu le malheur
d'aller à l'encontre du ``goût du jour'' et de ce fait semble impublia-%
ble, a néanmoins servi (et sert encore) comme un guide précieux dans
mes réflexions ultérieures, notamment dans le contexte de la géomé-%
trie algébrique absolue de caractéristique nulle.

Le cas où $(X,K)$ est une ``carte'' 2-dimensionnelle, i.e.
où les composantes connexes de $X\backslash K$ sont des 2-cellules ouvertes \unsure{.}
(et où de plus $K$ est muni d'un ensemble fini $S$ de ``sommets'',
tel que les composantes connexes de $K\backslash S$ soient des 1-cellules ou-%
vertes) attire progressivement mon attention dans les années suivantes.
La catégorie isotopique de ces cartes admet une description algébrique
particulièrement simple, via l'ensemble des ``repères'' (ou ``drapeaux'',
ou ``biarcs'') associés à la carte, qui se trouve naturellement muni
d'une structure d'ensemble à groupe d'opérateurs, sous le groupe
\[ \underline{C}_2 = \langle \sigma_0,\sigma_1,\sigma_2 \,|\, \sigma_0^{\,2} = \sigma_1^{\,2} = \sigma_2^{\,2} = (\sigma_0\sigma_2)^2 = 1 \rangle \ , \]
que j'appelle le \uline{groupe cartographique} (non orienté) de dimension 2.
Iladmet comme sous-groupe d'indice 2 le \uline{groupe cartographique orienté},
engendré par les produits en nombre pair des générateurs, qui peut
aussi se décrire comme
\[ \underline{C}_2^{+} = \langle \rho_s,\rho_f,\sigma \,|\, \rho_s\rho_f = \sigma,\ \sigma^2 = 1 \rangle \ , \]

\margin{12}

\finishscan
% END PAGE esq_119_012
% BEGIN PAGE esq_119_013
\scanpage{119/13}
\mspage{11}

(avec
\[ \rho_s = \sigma_2\sigma_1 \ , \quad \rho_f = \sigma_1\sigma_0 \ , \quad \sigma = \sigma_0\sigma_2 = \sigma_2\sigma_0 \ , \]
opérations de \uline{rotation élémentaire} d'un repère autour d'un sommet,
d'une face et d'une arête respectivement). Il y a un dictionnaire
parfait entre la situation topologique des cartes compactes, resp.
cartes compactes orientées, d'une part, et les ensembles finis à
groupe d'opérateurs $\underline{C}_2$ resp. $\underline{C}_2^+$ de l'autre, dictionnaire dont
l'existence était d'ailleurs plus ou moins connue, mais jamais
énoncée avec la précision nécessaire, ni développé tant soit peu.
Ce travail de fondements est fait avec le soin qu'il mérite dans
un excellent travail de DEA, fait en commun par Jean Malgoire et
Christine Voisin en 1976.

Cette réflexion prend soudain une dimension nouvelle,
avec cette remarque simple que le groupe $\underline{C}_2^+$ peut s'interpréter
comme un quotient du groupe fondamental d'une sphère orientée pri-
vée de trois points, numérotés 0, 1, 2, les opérations $\rho_s$, $\sigma$, $\rho_f$
s'interprétant comme les lacets autour de ces points, satisfaisant
la relation familière
\[ \ell_0\ell_1\ell_2 = 1 \ , \]
alors que la relation supplémentaire $\sigma^2 = 1$ i.e. $\ell_1^{\,2} = 1$ signi-%
fie qu'on s'intéresse au quotient du groupe fondamental correspon-%
dant à un indice de ramification imposé 2 au point 1, qui classi-%
fie donc les revêtements de la sphère, ramifiés au plus en les
points 0, 1, 2, avec une ramification égale à 1 ou 2 en les points
au-dessus de 1. Ainsi, les cartes orientées compactes forment une
catégorie isotopique équivalente à celle de ces revêtements, soumis
de plus à la condition supplémentaire d'être des revêtements finis.
Prenant maintenant comme sphère de référence la sphère de Riemann,
ou droite projective complexe, rigidifiée par les trois point 0, 1
et $\infty$ (ce dernier remplaçant donc 2), et se rappelant que tout revê-%
tement ramifié fini d'une courbe algébrique complexe hérite lui-même
d'une structure de courbe algébrique complexe, on aboutit à cette
constatation, qui huit ans après me paraît encore toujours aussi
extraordinaire : \uline{toute carte orientée} ``\uline{finie}'' \uline{se réalise canoniquement}
\uline{sur une courbe algébrique complexe} ! Mieux encore, comme la droite

\margin{13}

\finishscan
% END PAGE esq_119_013
% BEGIN PAGE esq_119_014
\scanpage{119/14}
\mspage{12}

projective complexe est définie sur le corps de base absolue $\mathbb{Q}$ ,
ainsi que les points de ramification admis, les courbes algébriques
obtenues sont définies non seulement sur $\mathbb{C}$, mais sur la clôture
algébrique $\overline{\mathbb{Q}}$ de $\mathbb{Q}$ dans $\mathbb{C}$. Quant à la carte de départ, elle se
retrouve sur la courbe algébrique, comme image inverse du segment
réel $[0,1]$ (où $0$ est considéré comme un sommet, et $1$ comme
milieu d'une ``arête pliée'' ayant 1 comme centre), lequel constitue
dans la sphère de Riemann la ``2-carte orientée universelle'' (*).
Les points de la courbe algébrique $X$ au-dessus de 0, de 1 et de $\infty$
ne sont autres que les sommets, et les ``centres'' des arêtes et des
faces respectivement de la carte $(X,K)$, et les ordres des sommets
et des faces ne sont autres que les multiplicités des zéros et
des pôles de la fonction rationnelle (définie sur $\overline{\mathbb{Q}}$) sur $X$,
exprimant sa projection structurale vers $\mathbb{P}^1_{\mathbb{C}}$.

Cette découverte, qui techniquement se réduit à si peu
de choses, a fait sur moi une impression très forte, et elle repré-%
sente un tournant décisif dans le cours de mes réflexions, un dé-%
placement notamment de mon centre d'intérêt en mathématique, qui
soudain s'est trouvé fortement localisé. Je ne crois pas qu'un
fait mathématique m'ait jamais autant frappé que celui-là, et ait
eu un impact psychologique comparable \notecall{copy1-2}{call-esq_119_014-2-1}{($^2$)}. Cela tient sûrement à la
nature tellement familière, non technique, des objets considérés,
dont tout dessin d'enfant griffonné sur un bout de papier (pour peu
que le graphisme soit d'un seul tenant) donne un exemple parfaitement
explicite. A un tel dessin se trouvent associés des invariants arith-%
métiques subtils, qui seront chamboulés complètement dès qu'on y
rajoute un trait de plus. S'agissant ici de cartes sphériques, don-%

\noindent\rule{4cm}{0.4pt}

(*) Il y a une description analogue des cartes finies non orientées,
éventuellement avec bord, en termes de courbes algébriques \uline{réelles},
plus précisément de revêtement de $\mathbb{P}^1_{\mathbb{R}}$ ramifié seulement on 0, 1, $\infty$,
la surface à bord associée à un tel revêtement étant $X(\mathbb{C})/\tau$, où $\tau$
est la conjugaison complexe. La carte non orientée ``universelle'' est
ici le disque, ou hémisphère supérieur de la sphère de Riemann, muni
comme précédemment du 1-complexe plongé $K = [0,1]$.

\margin{14}

\finishscan
% END PAGE esq_119_014
% BEGIN PAGE esq_119_015
\scanpage{119/15}
\mspage{13}

nant nécessairement naissance à des courbes de genre zéro (qui ne
fournissent donc pas des ``modules''), on peut dire que la courbe
en question est ``épinglée'' dès qu'on fixe trois de ses points, par
exemple trois sommets de la carte, ou plus généralement trois cen-%
tres de facettes (sommets, arêtes ou faces) - dès lors l'applica-%
tion structurale $f : X \to \mathbb{P}^{1}_{\mathbb{C}}$ peut s'interpréter comme une fonction
rationnelle
\[ f(z) = P(z)/Q(z) \in \mathbb{C}(z) \]
bien déterminée, quotient de deux polynômes bien déterminés pre-%
miers entre eux avec $Q$ unitaire, satisfaisant à des conditions
algébriques qui traduisent notamment le fait que $f$ soit non
ramifié en dehors des valeurs $0$, $1$, $\infty$, et qui impliquent que les
coefficients de ces polynômes sont des \uline{nombres algébriques} ; donc
leurs zéros sont des nombres algébriques, qui représentent res-%
pectivement les sommets et les centres des faces de la carte en-%
visagée.

Revenant au cas général, les cartes finies s'interpré-%
tant comme des revêtements sur $\bar{\mathbb{Q}}$ d'une courbe algébrique définie
sur le corps premier $\mathbb{Q}$ lui-même, il en résulte que le groupe de
Galois $\mathrm{I}\!\Gamma$ de $\bar{\mathbb{Q}}$ sur $\mathbb{Q}$ opère sur la catégorie de ces cartes de
façon naturelle. Par exemple, l'opération d'un automorphisme
$\gamma \in \mathrm{I}\!\Gamma$ sur une carte sphérique donnée par la fonction rationnelle
ci-dessus, est obtenue en appliquant $\gamma$ aux coefficients des po-%
lynômes $P$, $Q$. Voici donc ce mystérieux groupe $\mathrm{I}\!\Gamma$ intervenir comme
agent transformateur sur des formes topologico-combinatoires de la
nature la plus élémentaire qui soit, amenant à se poser des ques-%
tions comme : telles cartes orientées données sont-elles ``conjuguées'',
ou : quels exactement sont les conjuguées de telle carte orientée
donnée ? (Il y en a, visiblement, un nombre fini seulement.)

J'ai traité quelques cas concrets (pour des revêtements
de bas degrés) par des expédients divers, J. Malgoire en a traité
quelques autres - je doute qu'il y ait une méthode uniforme permet-%
tant d'y répondre à coups d'ordinateurs. Ma réflexion très vite

\margin{15}

\finishscan
% END PAGE esq_119_015
% BEGIN PAGE esq_119_016
\scanpage{119/16}
\mspage{14}

s'est engagée dans une directions plus conceptuelle, pour arriver
à appréhender la nature de cette action de $\mathrm{I\!\Gamma}$. On s'aperçoit
d'emblée que grosso-modo cette action est exprimée par une cer-%
taine action ``extérieure'' de $\mathrm{I\!\Gamma}$ sur le compactifié profini du
groupe cartographique orienté $\underline{C}_2^{+}$ , et cette action à son tour
est déduite par passage au quotient de l'action extérieure cano-%
nique de $\mathrm{I\!\Gamma}$ sur le groupe fondamental profini $\hat{\pi}_{0,3}$ de $(U_{0,3})_{\bar{\mathbb{Q}}}$ ,
où $U_{0,3}$ désigne la courbe-type de genre $0$ sur le corps premier
$\mathbb{Q}$, privée de trois points. C'est ainsi que mon attention s'est
portée vers ce que j'ai appelé depuis la ``\uline{géométrie algébrique}
\uline{anabélienne}'', dont le point de départ est justement une étude
(pour le moment limitée à la caractéristique zéro) de l'action
de groupes de Galois ``absolus'' (\unsure{nctamment|notamment} les groupes $\mathrm{Gal}(\bar{K}/K)$,
où $K$ est une extension de type fini du corps premier) sur des
groupes fondamentaux géométriques (profinis) de variétés algé-%
briques (définies sur $K$), et plus particulièrement (rompant avec
une tradition bien enracinée) des groupes fondamentaux qui sont
très éloignés des groupes abéliens (et que pour cette raison je
nomme ``\uline{anabéliens}''). Parmi ces groupes, et très proche du groupe
$\hat{\pi}_{0,3}$ , il y a le compactifié profini du groupe modulaire $S\ell(2,\mathbb{Z})$,
dont le quotient par le centre $\pm 1$ contient le précédent comme
sous-groupe de congruence mod 2, et peut s'interpréter d'ailleurs
également comme un groupe ``cartographique'' orienté, savoir celui
qui classifie les cartes orientées \uline{triangulées} (i.e. celles dont
les faces sont des triangles ou des monogones).

Toute carte finie orientée donne lieu à une courbe al-%
gébrique projective et lisse définie sur $\bar{\mathbb{Q}}$ , et il se pose alors
immédiatement la question : quelles sont les courbes algébriques
sur $\bar{\mathbb{Q}}$ obtenues ainsi - les obtiendrait-on toutes, qui sait ?
En termes plus savants, serait-il vrai que toute courbe algébrique
projective et lisse définie sur un corps de nombres interviendrait
comme une ``\unsure{ccurbe|courbe} modulaire'' possible pour paramétriser les cour-%
bes elliptiques munies d'une rigidification convenable ? Une telle
supposition avait l'air à tel point dingue que j'étais presque

\margin{16}

\finishscan
% END PAGE esq_119_016
% BEGIN PAGE esq_119_017
\scanpage{119/17}
\mspage{15}

gêné de la soumettre aux compétences en la matière. Deligne consulté
trouvait la supposition dingueen effet, mais sans avoir un contre-%
exemple dans ses manches. Moins d'un an après, au Congrès Interna-%
tional de Helsinki, le mathématicien soviétique Bielyi annonce
justement ce résultat, avec une démonstration d'une simplicité dé-%
concertante tenant en deux petites pages d'une lettre de Deligne -
jamais sans doute un résultat profond et déroutant ne fut démontré
en si peu de lignes !

Sous la forme où l'énonce Bielyi, son résultat dit essen-%
tiellement que \uline{toute courbe algébrique définie sur un corps de nom-}%
\uline{bres peut s'obtenir comme revêtement de la droite projective ramifié}
\uline{seulement en les points} 0, 1, $\infty$. Ce résultat semble être passé plus
ou moins inaperçu. Pourtant, il m'apparaît d'une portée considéra-%
ble. Pour moi, son message essentiel a été qu'\uline{il y a une identité}
\uline{profonde entre la combinatoire des cartes finies d'une part}, \uline{et la}
\uline{géométrie des courbes algébriques définies sur des corps de nombres},
\uline{de l'autre}. Ce résultat profond, joint à l'interprétation algébrico-%
géométrique des cartes finies, ouvre la porte sur un monde nouveau,
inexploré - et à portée de main de tous, qui passent sans le voir.

C'est près de trois ans plus tard seulement, voyant que
décidément les vastes horizons qui s'ouvrent là ne faisaient rien
trésaillir en aucun de mes élèves, ni même chez aucun des trois ou
quatre collègues de haut vol auxquels j'ai eu l'occasion d'en parler
de façon circonstanciée, que je fais un premier voyage de prospec-%
tion de ce ``monde nouveau'', de janvier à juin 1981. Ce premier jet se
matérialise en un paquet de quelques 1 300 pages manuscrites, bap-%
tisées ``La Longue Marche à travers la théorie de Galois''. Il s'agit
avant tout d'un effort de compréhension des relations entre groupes
de Galois ``arithmétiques'' et groupes fondamentaux profinis ``géomé-%
triques''. Assez vite il s'oriente vers un travail de formulation
calculatoire de l'opération de $\mathrm{Gal}(\bar{\mathbb{Q}}/\mathbb{Q})$ sur $\hat{\pi}_{0,3}$, et dans un
stade ultérieur, sur le groupe légèrement plus gros $S\ell(2,\mathbb{Z})^{\wedge}$, qui
donne lieu à un formalisme plus élégant et plus efficace. C'est au
cours de ce travail aussi (mais développé dans des notes distinctes)

\margin{17}

\finishscan
% END PAGE esq_119_017
% BEGIN PAGE esq_119_018
\scanpage{119/18}
\mspage{16}

qu'apparaît le thème central de la géométrie algébrique anabélienne,
qui est de reconstituer certaines variétés $X$ dites ``anabéliennes''
sur un corps absolu $K$ à partir de leur groupe fondamental mixte,
extension de $\mathrm{Gal}(\bar{K}/K)$ par $\pi_1(X_{\bar{K}})$ ; c'est alors que se dégage la
``conjecture fondamentale de la géométrie algébrique anabélienne'',
proche des conjectures de Mordell et de Tate que vient de démontrer
Faltings\notecall{copy1-3}{call-esq_119_018-3-1}{($^3$)}. C'est là aussi que s'amorcent une première réflexion sur
les groupes de Teichmüller, et les premières intuitions sur la struc-%
ture multiple de la ``tour de Teichmüller'' - les multiplicités modu-%
laires ouvertes $M_{g,\nu}$ apparaissant par ailleurs comme les premiers
exemples importants, en dimension $> 1$, de variétés (ou plutôt, de
multiplicités) qui semblent bien mériter l'appellation ``anabélienne''.
Vers la fin de cette période de réflexion, celle-ci m'apparaît com-%
me une réflexion fondamentale sur une théorie alors encore dans les
limbes, pour laquelle l'appellation ``\uline{théorie de Galois-Teichmüller}''
me semble plus appropriée que ``Théorie de Galois'' que j'avais d'abord
donnée à mes notes.

Ce n'est pas le lieu ici de donner un aperçu plus circons-%
tancié de cet ensemble de questions, intuitions, idées - y compris
des résultats palpables, certes. Le plus important me semble celui
signalé en passant au par. 2, savoir la fidélité de l'action exté-%
rieure de $\mathrm{I}\!\Gamma = \mathrm{Gal}(\bar{\mathbb{Q}}/\mathbb{Q})$ (et de ses sous-groupes ouverts) sur $\hat{\pi}_{0,3}$,
et plus généralement (si je me rappelle bien) sur le groupe fonda-%
mental de toute courbe algébrique ``anabélienne'' (i.e. dont le genre
$g$ et le ``nombre de trous'' $\nu$ satisfont l'inégalité $2g + \nu \geq 3$,
i.e. telle que $\chi(X) < 0$) définie sur une extension finie de $\mathbb{Q}$. Ce
résultat peut être considéré comme essentiellement équivalent au
théorème de Bielyi - c'est la première manifestation concrète, par
un énoncé mathématique précis, du ``message'' dont il a été question
plus haut.

Je voudrais terminer cet aperçu rapide par quelques mots
de commentaire sur la richesse vraiment inimaginable d'un groupe
anabélien typique comme le groupe $S\ell(2,\mathbb{Z})$ - sans doute le groupe
discret infini le plus remarquable qu'on ait rencontré, qui apparaît

\margin{18}

\finishscan
% END PAGE esq_119_018
% BEGIN PAGE esq_119_019
\scanpage{119/19}
\mspage{17}

sous une multiplicité d'avatars (dont certains ont été effleurés
dans le présent rapport), et qui du point de vue de la théorie de
Galois-Teichmüller peut être considéré comme la ``pierre de cons-%
truction'' fondamentale de la ``tour de Teichmüller''. L'élément de
structure de $S\ell(2,\mathbb{Z})$ qui me fascine avant tout, est bien sûr
l'action extérieure du groupe de Galois $\mathrm{I}\!\Gamma$ sur le compactifié
profini. Par le théorème de Bielyi, prenant les compactifiés pro-%
finis de sous-groupes d'indice fini de $S\ell(2,\mathbb{Z})$, et l'action ex-%
térieure induite (quitte à passer également à un sous-groupe ou-%
vert de $\mathrm{I}\!\Gamma$), \uline{on trouve essentiellement les groupes fondamentaux}
\uline{de toutes les courbes algébriques} (pas nécessairement compactes)
\uline{définis sur des corps de nombres} $K$, \uline{et l'action extérieure de}
$\mathrm{Gal}(\bar{K}/K)$ \uline{dessus} - du moins est-il vrai que tout tel groupe fon-%
damental apparaît comme quotient d'un des premiers groupes (*).
Tenant compte du ``yoga anabélien'' (qui reste conjectural), disant
qu'une courbe algébrique anabélienne sur un corps de nombres $K$
(extension finie de $\mathbb{Q}$) est connu à isomorphisme près quand on
connaît son groupe fondamental mixte (ou ce qui revient au même,
l'action extérieure de $\mathrm{Gal}(\bar{K}/K)$ sur son groupe fondamental profini
géométrique), on peut donc dire que \uline{toutes les courbes algébriques}
\uline{définies sur des corps de nombres sont} ``\uline{contenues}'' \uline{dans le compac-}%
\uline{tifié profini} $S\ell(2,\mathbb{Z})\hat{}$, \uline{et dans la connaissance d'un certain sous-}%
\uline{groupe} $\mathrm{I}\!\Gamma$ \uline{du groupe des automorphismes extérieurs de ce dernier} !
Passant aux abélianisés des groupes fondamentaux précédents, on voit
notamment que toutes les représentations abéliennes $\ell$-adiques chères
à Tate et consorts, définies par des jacobiennes et jacobiennes gé-%
néralisées de courbes algébriques définies sur des corps de nombres,
sont contenues dans cette seule action de $\mathrm{I}\!\Gamma$ sur le groupe profini
anabélien $S\ell(2,\mathbb{Z})\hat{}$!\notecall{copy1-4}{call-esq_119_019-4-1}{($^{4}$)}

\rule{5cm}{0.4pt}

(*) En fait, il s'agit de quotients de nature particulièrement tri-%
viale, par des sous-groupes abéliens produits de ``modules de Tate''
$\hat{\mathbb{Z}}(1)$, correspondants à des ``groupes-lacets'' autour de points à l'infini.

\margin{19}

\finishscan
% END PAGE esq_119_019
% BEGIN PAGE esq_119_020
\scanpage{119/20}
\mspage{18}

Il en est qui, face à cela, se contentent de hausser les
épaules d'un air  désabusé et de parier qu'il n'y a rien à tirer
de tout cela, sauf des rêves. Ils oublient, ou ignorent, que notre
science, et toute science, serait bien peu  de choses, si depuis
ses origines elle n'avait été nourrie des rêves et des visions
de ceux qui s'y adonnent avec passion.

4.        Dès le début de ma réflexion sur les cartes bidimension-%
nelles, je me suis intéressé plus particulièrement aux cartes dites
``\uline{régulières}'', c'est-à-dire celles dont le groupe des automorphismes
opère transitivement (et de ce fait, de façon simplement transitive)
sur l'ensemble des repères. Dans le cas orienté et en termes de
l'interprétation algébrico-géométrique du paragraphe précédent,
ce sont les cartes qui correspondent à un revêtement \uline{galoisien} de
la droite projective. Très vite aussi, et dès avant même qu'appa-%
raisse le lien avec la géométrie algébrique, il apparaît néces-%
saire aussi de ne pas exclure les cartes infinies, qui intervien-%
nent notamment de façon naturelle comme revêtements universels des
cartes finies. Il apparaît (comme conséquence immédiate du ``diction-%
naire'' des cartes, étendu au cas des cartes pas nécessairement finies)
que pour tout couple d'entiers naturels  $p, q$ $\geq$ 1, il existe à iso-%
morphisme (non unique) près une carte 1-\uline{connexe} et une seule qui
soit de type  $(p,q)$  i.e. dont tous les sommets soient d'ordre  $p$
et toutes les faces d'ordre  $q$, et cette carte est une carte régu-%
lière. Elle se trouve épinglée par le choix d'un repère, et son
groupe des automorphismes est alors canoniquement isomorphe au quo-%
tient du groupe cartographique (resp. du groupe cartographique
orienté, dans le cas orienté) par les relations supplémentaires
\[ \rho_s^{\,p} = \rho_f^{\,q} = 1. \]
Le cas où ce groupe est fini est le cas ``pythagoricien'' des cartes
régulières sphériques, le cas où il est infini donne les pavages

\margin{20}

\finishscan
% END PAGE esq_119_020
% BEGIN PAGE esq_119_021
\scanpage{119/21}
\mspage{19}

réguliers du plan euclidien ou du plan hyperbolique (*). Le lien
de la théorie combinatoire avec la théorie ``conforme'' des pavages
réguliers du plan hyperbolique était pressenti, avant qu'appa-%
raisse celui des cartes finies avec les revêtements finis de la
droite projective. Une fois ce lien compris, il devient évident
qu'il doit s'étendre également aux cartes infinies (régulières
ou non) : \uline{toute carte finie ou non, se réalise canoniquement sur
une surface conforme} (compacte si et seulement si la carte est
finie), \uline{en tant que revêtement ramifié de la droite projective com-%
plexe, ramifié seulement en les points} 0, 1, $\infty$. La seule difficulté
ici était de mettre au point le dictionnaire entre cartes topolo-%
giques et ensembles à opérateurs, qui posait quelques problèmes
conceptuels dans le cas infini, à commencer par la notion même de
``carte topologique''. Il apparaît nécessaire notamment, tant par
raison de cohérence interne du dictionnaire, que pour ne pas laisser
échapper certains cas intéressants de cartes infinies, de ne pas
exclure des sommets et des faces d'ordre infini. Ce travail de fon-%
dements a été fait également par J. Malgoire et C. Voisin, sur la
lancée de leur premier travail sur les cartes finies, et leur théo-%
rie fournit en effet tout ce qu'on était en droit d'en attendre
(et même plus...).

C'est en 1977 et 1978, parallèlement à deux cours de C 4
sur la géométrie du cube et sur celle de l'icosaèdre, que j'ai com-%
mencé à m'intéresser aux polyèdres réguliers, qui m'apparaissent
alors comme des ``réalisations géométriques'' particulièrement con-%
crètes de cartes combinatoires, les sommets, arêtes et faces étant
réalisés respectivement comme des points, des droites et des plans
dans un espace affine tridimensionnel convenable, avec respect des
relations d'incidence. Cette notion de réalisation géométrique

\rule{4cm}{0.4pt}

(*) Dans ces énoncés, il y a lieu de ne pas exclure le cas où $p$, $q$
peuvent prendre la valeur $+\infty$, qu'on rencontre notamment de façon
très naturelle comme pavages associés à certains polyèdres réguliers
infinis, cf. plus bas.

\margin{21}

\finishscan
% END PAGE esq_119_021
% BEGIN PAGE esq_119_022
\scanpage{119/22}
\mspage{20}

d'une carte combinatoire garde un sens sur un corps de base, et
même sur un anneau de base arbitraire. Elle garde également un
sens pour des polyèdres réguliers de dimension quelconque, en
remplaçant le groupe cartographique $\underline{C}_2$ par une variante n-dimen-%
sionnelle $\underline{C}_n$ convenable. Le cas n = 1, i.e. la théorie des
polygônes réguliers en caractéristique quelconque, fait l'objet
d'un cours DEA en 1977/78, et fait apparaître déjà quelques phé-%
nomènes nouveaux, comme aussi l'utilité de travailler non pas
dans un espace ambiant affine (ici le plan affine), mais dans
un espace \uline{projectif}. Ceci est dû notamment au fait que dans cer-%
taines caractéristiques (et notamment en caractéristique 2) le
centre d'un polyèdre régulier est rejeté à l'infini. D'autre
part, le contexte projectif, contrairement au contexte affine,
permet de développer avec aisance un formalisme de dualité pour
les polyèdres réguliers, correspondant au formalisme de dualité
des cartes combinatoires ou topologiques (où le rôle des sommets
et des faces, dans le cas n = 2 disons, se trouve interchangé).
Il se trouve que pour tout polyèdre régulier projectif, on peut
définir un hyperplan canonique associé, qui joue le rôle d'un
hyperplan à l'infini canonique, et permet de considérer le polyè-%
dre donné comme un polyèdre régulier affine.

L'extension de la théorie des polyèdres réguliers (et
plus généralement, de toutes sortes de configurations géométrico-%
combinatoires, y compris les systèmes de racines...) du corps de
base $\mathbb{R}$ ou $\mathbb{C}$ vers un anneau de base général, me semble d'une
portée comparable, dans cette partie de la géométrie, à l'exten-%
sion analogue qui a eu lieu depuis le début du siècle en géométrie
algébrique, ou depuis une vingtaine d'années en topologie ( *),
avec l'introduction du langage des schémas et celui des topos.

\rule{4cm}{0.4pt}

(*) En écrivant cela, je suis conscient que rares sont les topo-%
logues, encore aujourd'hui, qui se rendent compte de cet élargisse-%
ment conceptuel et technique de la topologie, et des ressources
qu'elle offre.

\margin{22}

\finishscan
% END PAGE esq_119_022
\end{document}
